\chapter{Objects, Arrows, and the Determinism Discipline}\label{ch03:chap}

This chapter defines the kernel's two ground-level notions --- objects and
arrows --- and develops the discipline that makes everything built on them
replayable. The mathematical content is thin by design: an object is a named,
explicitly ordered finite carrier, and an arrow is a total function between
two declared objects, equal to another arrow exactly when the declared data
agree. The engineering content is thick, and it is the chapter's real
subject. Three decisions do the work: object identity includes the declared
name, so that semantically distinct objects can never be silently coerced;
every order-sensitive constructor input must be a replayable ordered
sequence, with sets, mappings, and one-shot iterables rejected at
construction; and arrow equality is extensional over declared data only, with
presentation names excluded from equality and hashing. Each decision is
argued, not merely announced, and each is enforced in the constructors of the
reference implementation (\texttt{category.py} in \texttt{src/tks\_v8\_core})
and demonstrated by negative fixtures --- tests whose entire assertion is
that a forbidden construction fails with the right typed error. The chapter
closes with the deepest of the three arguments: a treatment of hash-order
nondeterminism as a general software-evidence problem, of which the kernel's
ordered-input rule is a local solution.

Throughout, statements about what the implementation does are statements
about the Gate~1 prototype sources hashed in the project manifest; their
evidential status is \textsc{executable validation} (ledger items
PO-TKS8-CORE-02 and PO-TKS8-CORE-03), which is a claim about the finite
executable model and its test suite, not a proof about all executions of all
Python interpreters.

\section{Objects}\label{ch03:sec:objects}

\subsection{Definition}

\begin{definition}[Finite object]\label{ch03:def:object}
An \emph{object} is a pair
\[
X = (\mathrm{name}_X,\ \mathrm{carrier}_X)
\]
where $\mathrm{name}_X$ is a non-blank string and $\mathrm{carrier}_X$ is a
finite, explicitly ordered sequence $(x_0, x_1, \dots, x_{n-1})$ of pairwise
distinct, hashable elements. We write $|X| = n$ for the carrier length. The
position of an element in the carrier is its \emph{standard index}; standard
indices run from $0$ to $|X| - 1$.

Two objects are \emph{equal} iff their names are equal \emph{and} their
carrier sequences are equal element-by-element in order.
\end{definition}

The implementation realizes Definition~\ref{ch03:def:object} as the frozen
dataclass \texttt{FiniteObject} with fields \texttt{name} and
\texttt{elements}. Construction validates all four conditions --- non-blank
string name, ordered-sequence input (Section~\ref{ch03:sec:ordered-rule}),
hashability, pairwise distinctness --- and raises the typed error
\texttt{ValidationError} on any violation. The class exposes
\texttt{index(element)}, returning the standard index of a carrier element
and raising \texttt{ValidationError} for a non-element; this method is the
single bridge between element-level talk and index-level talk, and
everything numerical downstream (mapping tuples, basis vectors, matrix
entries) lives on the index side.

Three consequences of the definition deserve immediate notice.

First, \emph{the carrier is a sequence, not a set}. The elements are
required to be distinct, so the carrier determines an underlying set, but
the object is strictly more data: the same three elements in two different
orders yield two \emph{unequal} objects. This is deliberate. The standard
indices are the coordinates against which arrows' mapping tuples, the
Hilbert-space basis of Chapter~7, and the shortlex machinery of Chapter~8
are all defined; an object that forgot its order would make all of them
ill-defined (Section~\ref{ch03:sec:ordered-rule}).

Second, \emph{hashability of elements is required but hash values are never
consulted for order}. Elements must be hashable so that fiber computations,
deduplication tables, and denotation keys can use them in dictionaries ---
as keys whose \emph{lookup} is hash-based but whose \emph{enumeration} the
kernel never performs. The distinction between using a hash table for
membership and using it for iteration is exactly the line drawn in
Section~\ref{ch03:sec:hash-order}.

Third, \emph{the name is load-bearing}. It is not a comment. It
participates in object equality, and through object equality in composition
legality, endomap constructibility, path well-typedness, and every
downstream judgment. The next section argues that this is the right choice.

\section{Name-inclusive identity}\label{ch03:sec:name-identity}

\subsection{The anti-coercion argument}

Consider the alternative: \emph{carrier-only identity}, under which two
objects are equal whenever their ordered carriers are equal. This is the
natural default --- it is what ``a finite set with an enumeration'' means
mathematically --- and it is wrong for a typed kernel, for one reason:
composition legality must track \emph{intent}, and carriers systematically
underdetermine intent.

The failure pattern is silent coercion between semantically distinct objects
of coincidentally equal carrier. Three escalating examples:

\begin{example}[Two registers over the same alphabet]\label{ch03:ex:two-registers}
A model has two binary registers, \texttt{GateA} and \texttt{GateB}, each
carrying the ordered carrier $(0, 1)$. Under carrier-only identity the two
objects are \emph{the same object}, and the kernel would accept the
composite of an arrow into \texttt{GateA} with an arrow out of
\texttt{GateB} --- a data-flow that the modeler never asserted exists.
Under name-inclusive identity the two objects are distinct, the composite
raises \texttt{CompositionTypeError}, and connecting the registers requires
writing an explicit arrow $\texttt{GateA} \to \texttt{GateB}$ --- an
auditable modeling decision, made once, in one reviewable place.
\end{example}

\begin{example}[Coincident carriers across pipelines]\label{ch03:ex:pipelines}
The scenario of Section~\ref{ch01:sec:untyped-scenario} ended with a
four-element evaluative carrier and a four-element quadrant carrier. Suppose
a refactor renames the quadrant labels so that both carriers become
literally $(e_0, e_1, e_2, e_3)$. Carrier-only identity now merges the two
objects, and the underived composite of that scenario --- the one that
smuggled a quadrant row into an evaluative conclusion --- becomes legal
again, silently, as a side effect of a renaming that touched no arrow and no
analysis. Name-inclusive identity is immune: \texttt{EvalSummary} and
\texttt{QuadState} remain distinct however their elements are spelled, and
the only way to re-enable the composite is to merge the declared object
identities themselves --- a change no reviewer would mistake for cosmetic.
\end{example}

\begin{example}[Same mathematics, different units]\label{ch03:ex:units}
Two objects carry $(0, 1, 2, 3)$: one is a count of retries, one an index
into a rubric. Every mathematical property of the two carriers is
identical. No inspection of elements, cardinalities, or even attached
arrows need reveal the difference --- the difference \emph{is} the intent,
and the name is the only place intent lives. Dimensional-analysis
disciplines in physical computation make the same move for the same reason:
quantities with identical numeric representations are kept apart by a
nominal tag because the cost of silent unit coercion is unbounded.
\end{example}

\begin{ruleenv}[Name-inclusive identity]\label{ch03:rule:name-identity}
Object equality is equality of the pair (declared name, ordered carrier).
No operation of the kernel --- composition, endomap construction, path
validation, lift construction, catalog registration --- ever identifies two
objects with different names, or two objects with the same name and
different carriers. There is no coercion path, implicit or explicit,
between distinct objects; the only way to relate them is an arrow.
\end{ruleenv}

The second half of the rule --- same name, different carriers, still
unequal --- closes the converse hole: a name is not a \emph{claim} that
overrides data, it is one component of a pair. Two versions of ``the same''
object with different carriers (an element added, an order changed) are
different objects, and arrows attached to the old version do not type-check
against the new. This is again deliberate: a carrier change alters standard
indices, hence the meaning of every mapping tuple that references the
object, and forcing the type mismatch to surface is strictly safer than
letting stale arrows flow into new coordinates.

\begin{remark}[Cost of the rule, and why it is paid]\label{ch03:rem:cost}
Name-inclusive identity has a real cost: it refuses some identifications a
mathematician would make without thought, and it obliges the modeler to
write explicit conversion arrows (often identities-on-carrier with a new
name) at boundaries where two vocabularies genuinely coincide. The kernel
pays this cost knowingly. The conversion arrow is not bureaucracy; it is
the audit record of the identification --- who asserted it, in which
version, reviewable and revocable. Every silent coercion the rule forbids
is a place where that record would otherwise not exist.
\end{remark}

\subsection{Carrier evolution and object versioning}

The same-name-different-carrier clause has a practical corollary worth
developing: it prescribes how models \emph{evolve}. Suppose an analysis
adds a fourth state to a three-state carrier. The tempting move ---
mutate ``the'' object in place --- is unavailable twice over: the
implementation's objects are frozen values, and the discipline's objects
are identity-carrying declarations. The correct move is to declare a new
object (in practice with a versioned name, say
\texttt{States\_v2}, though the kernel imposes no naming scheme) and to
relate it to the old one by explicit arrows: an inclusion
$\texttt{States} \to \texttt{States\_v2}$ embedding the old states, and,
where meaningful, a retraction collapsing the new state onto an old one.
Every existing arrow remains valid against the old object --- history
does not retype --- and participates in the new regime only through the
declared bridges, which are ordinary arrows subject to ordinary review.

The alternative regimes are both worse in characteristic ways. Silent
in-place mutation retypes history: mapping tuples written against the
old indices are reinterpreted against the new carrier, which is
Scenario~I of Chapter~\ref{ch01:chap} induced deliberately. And
carrier-only identity with mutation tracking (``the object is whatever
the carrier currently is'') makes object identity a function of time,
destroying the replayability of any digest computed before the change.
Name-inclusive identity with explicit versioning is the unique regime in
which every artifact ever produced remains exactly reinterpretable ---
its objects are immutable values that still mean what they meant. The
cost, once again, is explicitness; the payment, once again, is an audit
trail.

\subsection{Relation to the categorical picture}

Nothing in Rule~\ref{ch03:rule:name-identity} conflicts with the
mathematics of Chapter~\ref{ch02:chap}; it refines it. The kernel's
category (Chapter~4 proves it is one) simply has \emph{more objects} than
the category of bare finite carriers: one object per declared (name,
carrier) pair, with hom-sets empty between objects the modeler has not
connected. Categorically, names implement a discrete indexing of otherwise
isomorphic objects --- the standard device of working in a category where
isomorphic objects are deliberately not identified, because the
identification is exactly what must be controlled \cite{maclanecwm}. The
mathematical content of any construction is unchanged by the extra tags;
the \emph{legality} of compositions is what the tags govern.

\section{The ordered-input rule}\label{ch03:sec:ordered-rule}

\subsection{Statement and enforcement point}

\begin{ruleenv}[Ordered input]\label{ch03:rule:ordered-input}
Every constructor input whose order affects any downstream artifact must be
supplied as a replayable ordered sequence. Sets, mappings, and one-shot
iterables are rejected at construction time with \texttt{ValidationError}.
The rule is enforced for exactly seven input positions:
\begin{enumerate}
\item carrier elements (of a \texttt{FiniteObject});
\item arrow mappings (the image-index tuple of a \texttt{FiniteArrow});
\item path arrows (the arrow sequence of a \texttt{SemigroupoidPath});
\item word labels (the generator-label sequence of a typed word);
\item catalog objects (the object sequence of a generator catalog);
\item catalog generators (the generator sequence of a catalog, whose
declared order is the shortlex ranking);
\item lift coordinate vectors (coefficient sequences fed to the
deterministic lift).
\end{enumerate}
\end{ruleenv}

The enforcement point is a single shared validator,
\texttt{\_ordered\_tuple}, called by every constructor listed. Its
acceptance condition is deliberately conservative: the input must be a
sequence in the collections sense --- indexable, re-iterable, with an
order that is part of its contract --- and must not be one of the
explicitly excluded types. Excluded are:

\begin{itemize}
\item \emph{sets} (any \texttt{Set} instance): iteration order is not part
of the contract and, for hash-randomized element types, varies by process
(Section~\ref{ch03:sec:hash-order});
\item \emph{mappings} (any \texttt{Mapping} instance): even where a
concrete mapping type happens to iterate deterministically, the abstract
contract does not promise an order, and accepting mappings would invite
callers to encode order-sensitive data in a container whose ordering
guarantees vary by type, language version, and implementation;
\item \emph{one-shot iterables} (generators, iterators, streams): a
consumed iterator re-iterates as empty, so validation and storage would
observe different data than a second traversal; nothing about the values
is stable enough to constitute a declaration;
\item \emph{strings and byte strings} (\texttt{str}, \texttt{bytes},
\texttt{bytearray}): these \emph{are} ordered sequences, and they are
rejected anyway, because a string supplied where a sequence of elements is
expected almost always signals an error --- it would silently explode into
a carrier of single characters, an object the caller never intended.
Rejecting the container forces the caller to write the character sequence
explicitly if that is truly what is meant.
\end{itemize}

What remains --- lists, tuples, and other genuine sequences --- is
converted to an immutable tuple at the constructor boundary, so that the
stored order can never drift after validation.

\subsection{Why exactly these seven positions}

The rule is not a blanket ban on unordered containers; the kernel freely
uses sets and dictionaries \emph{internally} where order does not reach any
output --- membership tests, duplicate detection, fiber grouping keyed by
element. The seven positions of Rule~\ref{ch03:rule:ordered-input} are
exactly the places where order is \emph{semantics}:

\emph{Carrier elements} fix the standard indices, and standard indices are
the coordinate system for everything: an arrow's mapping tuple is read
against them (position $i$ of the tuple is the image of carrier element
$i$), and the Hilbert space $H_X$ of Chapter~7 takes its orthonormal basis
$(e_{x_0}, \dots, e_{x_{n-1}})$ in carrier order, so that permuting the
carrier conjugates every lift matrix by a permutation. Two enumerations of
the same element set yield different objects with different geometry ---
correctly so, and only tolerably so if the enumeration was chosen by the
caller rather than by a hash table.

\emph{Arrow mappings} are the function itself: the tuple
$(t_0, \dots, t_{n-1})$ \emph{is} the total function in the tuple
representation of Section~\ref{ch02:sec:finsets}. An unordered or
re-orderable source for this data is an unordered function --- a
contradiction.

\emph{Path arrows} fix which composite a path denotes; composition is not
commutative, and a path delivered as a set of arrows would denote as many
composites as orderings.

\emph{Word labels} fix both the denotation of a typed word (the composite
of its arrows in sequence) and its identity as a word --- the thing
shortlex compares letter by letter.

\emph{Catalog objects and catalog generators} fix the closure's universe
and, critically, the \emph{ranking} $\rho$ of
Definition~\ref{ch02:def:shortlex}: the declared generator order is the
lexicographic tiebreak of shortlex, so the canonical representative of
every denotation depends on it
(Remark~\ref{ch02:rem:slex-typed}). A catalog fed from a set would have a
per-process canonical form --- the contradiction in terms already noted in
Section~\ref{ch01:sec:nondet-scenario}.

\emph{Lift coordinate vectors} pair coefficients with basis vectors by
position; an unordered coefficient container does not determine a vector.

In every one of the seven, the downstream artifact is either an equality
judgment, a digest, or a canonical form. That is the criterion. Order that
reaches none of these may be, and is, left to internal convenience.

\subsection{Order sensitivity made concrete: a two-element worked example}

The abstract claim ``carrier order changes downstream artifacts'' can be
made completely concrete at the smallest nontrivial size, and the
exercise is worth a page because it shows the change reaching all three
artifact classes --- equality, matrices, and canonical forms --- from a
single transposition.

Let the element set be $\{u, v\}$ and consider the two enumerations:
$X = (\texttt{Pair}, (u, v))$ and $X' = (\texttt{Pair}, (v, u))$. These
are unequal objects (same name, different carrier order), and the
intended map ``send everything to $u$'' is a different tuple over each:
over $X$ it is the arrow $c$ with mapping $(0, 0)$ (both positions to
index $0$, which is $u$); over $X'$ it is the arrow $c'$ with mapping
$(1, 1)$ (both positions to index $1$, which is $u$ there).

\emph{Equality:} $c$ and $c'$ are extensionally unequal --- different
sources, different tuples --- although both realize ``constantly $u$''
on the same element set. Any digest over denotation keys distinguishes
them; a system that conflated them would be conflating the coordinate
systems in which every other arrow over these objects is written.

\emph{Matrices:} the lift of $c$ over $X$ has matrix rows
$(1, 1)$ and $(0, 0)$ --- column $j$ has its $1$ in row $0$ because
both elements map to index $0$ --- while the lift of $c'$ over $X'$ has
rows $(0, 0)$ and $(1, 1)$. The two matrices are permutation-conjugate
(Lemma~\ref{ch02:lem:perm-similarity} territory: same spectrum-relevant
structure, same norm $\sqrt{2}$), but they are \emph{different matrices},
and any artifact that serializes matrix entries --- a test expectation, a
digest, a Lean fixture --- sees the difference. Spectral invariants
survive reordering; byte-level evidence does not.

\emph{Canonical forms:} put both arrows in generator catalogs. Over
$X$, a catalog declaring generators in the order
$(c, \id_X)$ ranks $c$ as letter $0$; the same semantic content
delivered in the other declared order ranks it as letter $1$, and every
shortlex tiebreak between words of equal length can flip
(Example~\ref{ch02:ex:slex-two} showed the flip at alphabet level).
Carrier order and catalog order are thus two independent dials, each of
which the kernel forces into explicit, versioned data.

None of this is pathology; it is coordinates behaving as coordinates.
The pathology would be letting a hash table set the dials.

\begin{invariant}[Replayability of constructed values]\label{ch03:inv:replay}
Every value constructible through the kernel's public constructors is a
pure function of its constructor arguments read as ordered data. Two calls
with equal ordered arguments, in any two processes on any two machines,
construct equal values with equal serializations. No kernel output depends
on hash seeds, iteration order of unordered containers, object identity
(memory addresses), wall-clock time, or any other ambient interpreter
state.
\end{invariant}

Invariant~\ref{ch03:inv:replay} is the chapter's central promise, and
Rule~\ref{ch03:rule:ordered-input} is most of its enforcement. The
remaining exposure --- the caller might \emph{produce} the ordered argument
nondeterministically before the kernel ever sees it --- cannot be closed by
a library and is addressed honestly: the kernel's constructors make the
ordering decision the caller's explicit, visible act, which is the most a
constructor can do. Section~\ref{ch03:sec:hash-order} explains why so much
weight is placed here.

\section{Hash-order nondeterminism as an evidence problem}\label{ch03:sec:hash-order}

This section develops, at mechanism level, the failure that
Rule~\ref{ch03:rule:ordered-input} exists to exclude. Although the concrete
details are Python's, the problem is general, and we state it generally
before and after the specifics.

\subsection{The mechanism}

Hash-based containers --- sets and hash maps --- store elements in an
internal table at positions derived from the elements' hash values.
Iteration walks the table, so iteration order is a function of (i) the hash
values of the elements, (ii) the table's size history, which depends on
insertion order and interim deletions, and (iii) implementation details of
collision resolution. None of these is part of the container's semantic
contract; the contract is membership.

In CPython specifically, the hash of a string is additionally
\emph{randomized per process}: since Python 3.3, string and bytes hashing
incorporates a random key generated at interpreter startup, as a hardening
measure against algorithmic-complexity denial-of-service attacks on hash
tables (the modern implementations use keyed hashes of the SipHash
family). The environment variable \texttt{PYTHONHASHSEED} controls the
key: unset, the key is random per process; set to an integer, the key is
that integer (fixed across processes); set to \texttt{0}, randomization is
disabled entirely. Consequently:

\begin{itemize}
\item the set $\{ \texttt{"calm"}, \texttt{"alert"}, \texttt{"overload"}
\}$ iterates in an order that varies from process to process under default
settings;
\item small integers hash to themselves, so sets of small integers
typically iterate in a stable-looking order --- which makes int-keyed
prototypes appear deterministic and lets the defect ship, to detonate later
when the element type changes to strings, tuples, or enums;
\item \texttt{dict} preserves insertion order as a language guarantee in
modern Python, so a mapping built by deterministic insertions does iterate
deterministically --- but a mapping built from set-derived data inherits
the upstream nondeterminism, and the abstract \texttt{Mapping} contract
promises nothing; the kernel rejects the whole class rather than
adjudicating pedigrees it cannot see.
\end{itemize}

\subsection{How the nondeterminism reaches digests}

The leak path is short and typical. Order-sensitive data is assembled by
iterating a hash-ordered container --- \texttt{tuple(states)},
\texttt{list(mapping)}, a comprehension over a set, a \texttt{join} over
\texttt{dict.keys()} of a set-fed dict. The assembled sequence is then
serialized --- into a JSON document, a repr string, a pickle, a canonical
byte encoding --- and the serialization is digested with a cryptographic
hash (in this project, SHA-256) whose output is recorded as evidence: in a
manifest, a test log, a signed receipt, a changelog entry. Every step
after the first is deterministic; the digest is nonetheless a random
variable, because its input embedded an iteration order.

The scenario of Section~\ref{ch01:sec:nondet-scenario} traced one such
path end to end: a three-element set becomes a carrier, the carrier order
fixes standard indices, an intended-constant endomap serializes as
$(1,1,1)$ in one process and $(2,2,2)$ in another, and the recorded digest
fails to reproduce under the auditor's replay. Note the exact epistemic
damage: the digest does not become \emph{weaker} evidence --- it becomes
\emph{no} evidence, because a verifier that fails on honest replays cannot
distinguish honest divergence from tampering or corruption. And the
failure is maximally confusing to diagnose, because every individual
function in the pipeline is correct in isolation; the defect lives in a
container choice that type-checks, lints clean, and passes every
single-process test.

\begin{threat}[Hash-order leakage]\label{ch03:threat:hash-order}
\emph{Asset}: replayability of every digest, canonical form, and equality
judgment produced by the kernel and recorded in its manifests.
\emph{Vector}: any construction of order-sensitive data by iteration over
a hash-ordered or otherwise contract-free container, anywhere upstream of
a recorded artifact --- including in caller code, test fixtures, and data
preparation scripts. \emph{Amplifiers}: element types whose hashes are
process-seeded (strings, bytes, composites containing them);
deterministic-looking prototypes over small integers; mappings of
unverifiable pedigree; iterator consumption that silently yields different
data to validation and storage. \emph{Kernel mitigation}:
Rule~\ref{ch03:rule:ordered-input}, enforced at all seven order-sensitive
constructor positions, converting the vector into an immediate typed
error. \emph{Residual exposure}: nondeterministic assembly in caller code
prior to the constructor boundary; addressed by making the ordering an
explicit argument and by the manifest practice of digesting inputs as well
as outputs, so a nondeterministic caller fails to reproduce its own inputs
and is caught at the evidence layer.
\end{threat}

\subsection{A minimal reproduction experiment}

The mechanism can be exhibited in two interpreter invocations, and we
record the experiment in a form any reader can replay, because the
failure's invisibility in source code is best dispelled by watching it
happen. Run the same one-liner twice as separate processes:

\begin{quote}
\texttt{python -c "print(tuple(\{'calm', 'alert', 'overload'\}))"}
\end{quote}

Under default settings (hash randomization on), the two processes will,
with high probability across a handful of repetitions, print the three
strings in different orders --- same program text, same input data,
different output. Fixing the seed makes the accident repeatable without
making it principled:

\begin{quote}
\texttt{PYTHONHASHSEED=0 python -c "print(tuple(\{'calm', 'alert',
'overload'\}))"}
\end{quote}

now prints the same order on every invocation --- \emph{some} order,
chosen by the interpreter's hash function at that version, promised by
nothing, and subject to change when the interpreter changes. Repeating
the experiment with the integer set \texttt{\{1, 2, 3\}} typically
prints \texttt{(1, 2, 3)} stably under any seed --- the deceptive
stability of small-integer hashing noted above, and the reason
prototypes over integers ship the defect. The corresponding kernel
experiment is the negative fixture: constructing a \texttt{FiniteObject}
from any of these set literals raises \texttt{ValidationError} under
every seed, which is the only outcome of the four that is both
deterministic and honest.

For contrast, the deterministic path through the kernel:

\begin{quote}
\texttt{obj = FiniteObject(name="States", elements=("alert", "calm",
"overload"))}\\
\texttt{arrow = FiniteArrow(name="to-calm", source=obj, target=obj,
mapping=(1, 1, 1))}
\end{quote}

constructs, in every process on every machine, an object whose element
\texttt{"calm"} has standard index $1$ and an arrow whose serialized
denotation key is byte-identical across runs --- the property the
digest discipline below consumes.

\subsection{Digest discipline at the evidence layer}

The kernel's outputs become evidence when they are serialized and
digested, so the discipline extends one level beyond the constructors.
Three rules govern the project's digest practice; they are stated here
because they complete the determinism story, though their enforcement
lives in the manifest tooling rather than in \texttt{category.py}.

\emph{Digest denotations, not presentations.} Semantic digests are
computed over denotation keys --- source object (name and ordered
carrier), target object, mapping tuple --- and never over presentation
names, memory addresses, or repr strings of container types whose
formatting is version-dependent. Two arrows that are extensionally
equal must digest identically, or the digest would distinguish what the
mathematics identifies.

\emph{Serialize with a fixed, versioned encoding.} The byte encoding of
a denotation key (field order, integer encoding, string encoding,
delimiter convention) is itself part of the evidence contract: it is
fixed, documented, and versioned, so that a digest names one byte
string under one encoding version. An ``obvious'' serialization left to
a library default is a hash-order problem one level up --- defaults
change across library versions exactly as iteration orders change
across interpreter versions.

\emph{Digest inputs alongside outputs.} A manifest records digests of
the inputs to a computation as well as its outputs, so that a
nondeterministic producer upstream of the kernel is caught by failing to
reproduce its own inputs, rather than remaining invisible behind a
reproducible kernel. This is the residual-exposure mitigation named in
Threat~\ref{ch03:threat:hash-order}: the kernel cannot make caller code
deterministic, but the evidence layer can make caller nondeterminism
\emph{visible}.

\subsection{The general software-evidence problem}

Nothing above is specific to Python, to this kernel, or even to hashing.
The general form: \emph{a computation is offered as evidence exactly to the
extent that an independent party can replay it to the same bytes; any
dependence on ambient state --- hash seeds, iteration order, thread
interleavings, locale, filesystem enumeration order, clock, address-space
layout --- silently converts a replayable computation into an unreplayable
one, and the conversion is invisible at the point of use.} The
reproducible-builds discipline in software distribution fights precisely
this class (embedded timestamps, archive member order, parallel link
order); numerical computing fights its floating-point cousin (reduction
order under parallelism); and every attestation scheme that hashes an
artifact inherits every nondeterminism of the artifact's producer. The
principled responses available to a system designer are three: (a)
\emph{canonicalize at the boundary} --- sort, normalize, and fix an order
before anything order-sensitive is derived; (b) \emph{refuse at the
boundary} --- reject inputs whose order is not part of their contract, and
force the caller to decide; (c) \emph{pin the ambient state} --- fix seeds
and environments so that the accident is at least a repeatable accident.

The kernel chooses (b), and it is worth recording why the two alternatives
were declined. Canonicalization (a) requires a canonical order on
arbitrary hashable elements, which does not exist in general (elements
need not be mutually comparable) and, where it does exist, encodes a
semantic choice --- \emph{which} order --- that belongs to the modeler, not
the library; a silently sorted carrier is exactly the kind of unreviewed
decision the kernel refuses to make. Pinning (c) --- running everything
under a fixed \texttt{PYTHONHASHSEED} --- makes runs repeatable but leaves
the semantics coupled to an environment variable: the ``canonical'' form
becomes a function of a deployment setting, unreproducible by any party
who does not know it, and fragile under interpreter upgrades that change
hash functions. Refusal (b) is the only option under which the constructed
value is a pure function of reviewable arguments
(Invariant~\ref{ch03:inv:replay}); its cost --- callers must write
\texttt{sorted(...)} or an explicit literal where they previously wrote a
set --- is paid once, visibly, at the point where the decision belongs.

\begin{boundary}\label{ch03:boundary:evidence}
The kernel's determinism guarantees are claims about the kernel's own
constructors and outputs (Invariant~\ref{ch03:inv:replay}), validated by
the negative fixtures of Section~\ref{ch03:sec:negative-fixtures} on the
tested platforms. They are not a claim that programs \emph{using} the
kernel are deterministic, nor a general theorem about Python semantics,
nor a certification of any digest produced outside the kernel's
constructor discipline.
\end{boundary}

\section{Arrow judgments and extensional equality}\label{ch03:sec:arrows}

\subsection{The arrow judgment}

\begin{definition}[Arrow]\label{ch03:def:arrow}
An \emph{arrow judgment} $f : X \to Y$ consists of a declared source
object $X$, a declared target object $Y$, and a total map from the
standard indices of $X$ to the standard indices of $Y$, represented as the
\emph{mapping tuple} $(f(0), f(1), \dots, f(|X|-1))$ with each entry in
$\{0, \dots, |Y|-1\}$.
\end{definition}

The implementation is the frozen dataclass \texttt{FiniteArrow} with
fields \texttt{name}, \texttt{source}, \texttt{target}, and
\texttt{mapping}. Construction validates: the name is a non-blank string;
source and target are \texttt{FiniteObject} values; the mapping passes the
ordered-input rule and has length exactly $|X|$ (totality is syntactic:
one image per source index, no more, no fewer); and every entry is an
integer within range for the target. Two validation details are easy to
miss and are load-bearing:

\emph{Booleans are rejected as indices.} In Python, \texttt{bool} is a
subclass of \texttt{int}, so \texttt{True} would silently function as
index $1$. The constructor tests for \texttt{bool} explicitly and rejects
it: an index that arrived as a truth value is almost certainly a logic
error upstream, and the coincidence \texttt{True} $= 1$ is exactly the
kind of accidental meaning the kernel exists to refuse.

\emph{Range checking is against the declared target, not against use.} An
out-of-range entry fails at construction, not at first application; there
is no constructible arrow whose application can index outside its target's
carrier.

Application is then defined elementwise: for a carrier element $x$ of $X$,
the arrow returns the target carrier element at position
$\mathrm{mapping}[\mathrm{index}(x)]$ --- element in, element out, with
the index representation confined to the interior.

\subsection{Element-level application and the index bridge}

Application deserves a worked micro-example, because it is where the
two vocabularies of the kernel --- elements, which callers care about,
and indices, which the mathematics runs on --- meet, and because the
meeting point is where a lesser design would leak coordinates to the
caller or elements into the algebra.

Take $X = (\texttt{Signal}, (\texttt{low}, \texttt{high}))$ and
$Y = (\texttt{Action}, (\texttt{hold}, \texttt{raise},
\texttt{drop}))$, and the arrow \texttt{policy} $: X \to Y$ with
mapping $(1, 2)$. Applying \texttt{policy} to the element
\texttt{"low"} proceeds in three steps, each total and each validated:
the source's \texttt{index} method resolves \texttt{"low"} to standard
index $0$ (raising \texttt{ValidationError} for any non-element, so
application cannot be fed a stray value); the mapping tuple sends $0$
to $1$; the target's carrier resolves index $1$ to the element
\texttt{"raise"}, which is returned. Elements in, elements out; the
integer $1$ existed only in the interior. The caller never handles
indices, so no caller can develop a dependency on them; the algebra
never handles elements, so no element's incidental properties (its
hash, its string form, its comparability) can influence a computation.
Every structure of Chapters~4 through~8 --- composition, powers, lifts,
words --- operates purely on the index side, with
\texttt{index} and carrier lookup as the only crossings, both owned by
the object that defines the coordinates.

The design also fixes where partiality is allowed to appear:
\emph{application to a non-element} is a caller error and raises;
\emph{application of a validated arrow to a carrier element} cannot
fail, by the range validation at construction. There is no code path
on which a constructed arrow returns an out-of-carrier value, and the
test suite's positive families exercise exactly this closure property
across every arrow they enumerate.

\subsection{Extensional equality}

\begin{definition}[Extensional equality of arrows]\label{ch03:def:ext-eq}
Arrows $f$ and $g$ are \emph{extensionally equal}, written $f = g$, iff
their sources are equal objects, their targets are equal objects, and
their mapping tuples are equal. The triple
$(\mathrm{source}, \mathrm{target}, \mathrm{mapping})$ is the arrow's
\emph{denotation key}; extensional equality is equality of denotation
keys.
\end{definition}

\begin{proposition}[Equality is a congruence]\label{ch03:prop:congruence}
Extensional equality is an equivalence relation on arrows, and composition
respects it: if $f = f'$ and $g = g'$ and $g \circ f$ is defined, then
$g' \circ f'$ is defined and $g \circ f = g' \circ f'$. Likewise
application respects it: $f = f'$ implies $f(x) = f'(x)$ for every
carrier element $x$ of the common source.
\end{proposition}

\begin{proof}
Equivalence is immediate: equality of triples of (name-inclusive) objects
and tuples is reflexive, symmetric, and transitive componentwise.
For the congruence: $f = f'$ gives equal targets and $g = g'$ equal
sources, so definedness of $g \circ f$ --- target of $f$ equals source of
$g$ as objects --- transfers to $g' \circ f'$. The composite's denotation
key is determined by the operands' keys: source of $f$, target of $g$,
and the tuple whose $i$-th entry is
$\mathrm{mapping}_g[\mathrm{mapping}_f[i]]$; equal operand keys therefore
force equal composite keys. Application reads only the source, mapping,
and target, all fixed by the key.
\end{proof}

\begin{corollary}[Derived operations respect equality]\label{ch03:cor:derived-congruence}
Extensional equality is respected by every derived operation of the
kernel: if $f = f'$ then $f^{n} = f'^{\,n}$ for every $n \geq 0$ (when
the endomap constructions are defined, which they are for both or
neither); if two paths have componentwise equal arrow sequences, their
composites are equal; and equal arrows produce equal lifts, equal
closure-table entries, and equal normal-form assignments in
Chapters~7 and~8.
\end{corollary}

\begin{proof}
Powers: induction on $n$ from Proposition~\ref{ch03:prop:congruence},
with base case $f^{0} = \id_X = f'^{\,0}$ (equal sources give equal
identity arrows). Definedness transfers because \texttt{is\_endomap} is
a predicate on the denotation key (source equals target as objects).
Path composites: induction on path length using the congruence at each
fold step. The lift and closure statements are established in their own
chapters; each construction there reads only denotation keys, which is
all the present corollary requires of them.
\end{proof}

Proposition~\ref{ch03:prop:congruence} and its corollary are short but
foundational: they say that every construction of the kernel factors
through denotation keys, so that reasoning may be conducted entirely on
the extensional data. In the implementation, \texttt{denotation\_key} is
exposed as a property precisely so that higher layers --- the
deduplication table of the normal-form closure in Chapter~8 above all
--- can key on the mathematical identity of an arrow rather than on any
incidental attribute. Which brings us to the one incidental attribute an
arrow carries.

\section{Presentation names are not identity}\label{ch03:sec:presentation-names}

Every arrow has a \texttt{name}: a non-blank string, required at
construction, and used in every diagnostic the kernel emits ---
composition mismatch errors print
\texttt{f:X->Y} then \texttt{g:Z->W} in terms of these names, and the
error messages of Chapters~5 through~8 do likewise. Names are also
synthesized by the kernel itself: the identity constructor names its
result \texttt{id[X]} after the object; composition names its result
\texttt{(g o f)} after the operands; the power operation names
$f^{n}$ with a compact exponent label; the path composite is named by its
length. For very large exponents the synthesized name uses the exponent's
bit length plus a short SHA-256 label rather than full decimal rendering,
so that name construction is logarithmic in the exponent and immune to the
host language's integer-to-decimal conversion limits (Chapter~5 details
this; regression tests cover exponents $2^{30}$ and $10^{5000}$).

\begin{ruleenv}[Presentation names never enter equality]\label{ch03:rule:names-out}
The \texttt{name} field of an arrow is presentation data: it participates
in no equality comparison, no hash, no denotation key, no digest of
semantic content, and no decision of any kernel algorithm. Two arrows with
equal denotation keys are equal whatever their names; two arrows with
distinct keys are unequal whatever their names. In the implementation the
field is declared \texttt{compare=False}, excluding it from the generated
equality and hash of the frozen dataclass.
\end{ruleenv}

The rule resolves a tension that Section~\ref{ch03:sec:name-identity}
created deliberately. \emph{Object} names are semantic --- they carry
modeling intent that the carrier underdetermines, and so they participate
in identity. \emph{Arrow} names are not --- an arrow's full semantic
content is already present in its denotation key (source object including
its name, target object including its name, mapping tuple), and a name
over and above that content can only describe \emph{how the arrow was
obtained}: as a composite, as a power, as somebody's labeled generator.
History of construction is exactly what extensional mathematics discards
--- $f^{6}$ computed by binary powering and the same map entered as a
literal table are the same arrow --- and letting names into equality would
make the kernel intensional: normal forms would distinguish denotations by
their derivations, the closure's deduplication would fail to merge equal
maps, and the count of ``distinct'' arrows would depend on the syntactic
route taken to each. The two rules are therefore two sides of one
principle: \emph{identity includes exactly the data that composition
legality and denotation need --- all of it, and nothing else.} For
objects that includes the name; for arrows it excludes it.

One consequence is worth flagging for users: since names do not
participate in equality, they are also not \emph{protected} by it ---
nothing stops two extensionally different arrows from bearing the same
name, and diagnostics that print names can therefore be ambiguous where
keys are not. The kernel's diagnostics accept this, always printing
source and target names alongside, and the evidence layer never relies on
arrow names for anything.

\section{Negative-fixture testing of the constructors}\label{ch03:sec:negative-fixtures}

\subsection{Rejection behavior is specified behavior}

A kernel whose safety story is ``invalid things are unwritable'' owes its
users evidence that invalid things are, in fact, unwritable. Positive
tests cannot supply this: a suite that only exercises valid constructions
would pass unchanged against an implementation that accepted everything.
The kernel's suite therefore treats \emph{rejection} as first-class
specified behavior, tested by \emph{negative fixtures}: minimal invalid
inputs, each asserted to raise the exact typed error ---
\texttt{ValidationError} for malformed data,
\texttt{CompositionTypeError} for type mismatches --- at the exact
boundary, with no partial construction surviving. This is the executable
content of ledger items PO-TKS8-CORE-02 (reject every source/target
mismatch) and PO-TKS8-CORE-03 (keep one-carrier powers distinct from
heterogeneous paths), both of status \textsc{executable validation}.

\begin{checkclass}[Negative fixtures for object construction]\label{ch03:check:objects}
Each of the following must raise \texttt{ValidationError}:
a blank or non-string object name; carrier elements supplied as a
\texttt{set}, a \texttt{dict}, a generator expression, or a string;
carrier elements containing a duplicate; carrier elements containing an
unhashable value (e.g.\ a list). Additionally, \texttt{index} on a
non-element must raise \texttt{ValidationError} naming the object.
\end{checkclass}

\begin{checkclass}[Negative fixtures for arrow construction]\label{ch03:check:arrows}
Each of the following must raise \texttt{ValidationError}:
a blank arrow name; a source or target that is not a
\texttt{FiniteObject}; a mapping supplied as a set, mapping, or one-shot
iterable; a mapping tuple whose length differs from $|X|$ (both too short
--- a partial function --- and too long); a mapping entry that is a
\texttt{bool}; a mapping entry that is not an integer; a mapping entry
outside $\{0, \dots, |Y|-1\}$ (both negative and $\geq |Y|$).
\end{checkclass}

\begin{checkclass}[Negative fixtures for typed operations]\label{ch03:check:ops}
Each of the following must raise \texttt{CompositionTypeError}:
composing $g \circ f$ where the target of $f$ and source of $g$ differ in
name only (equal carriers); differ in carrier only (equal names); differ
in carrier order only (equal names, equal element sets); constructing an
\texttt{Endomap} from an arrow whose source and target are distinct
objects --- including the same-carrier different-name case; constructing
a \texttt{SemigroupoidPath} containing one adjacent mismatch of any of
the above kinds. Each of the following must raise
\texttt{ValidationError}: an empty path; a path containing a non-arrow;
a negative or non-integer or boolean endomap exponent.
\end{checkclass}

The three same-carrier variants inside
Checkclass~\ref{ch03:check:ops} are the anti-coercion argument of
Section~\ref{ch03:sec:name-identity} turned into executable assertions:
they pin exactly the coercions --- by name, by carrier, by order --- that
a weaker identity discipline would silently perform, and they fail
against any implementation that relaxes
Rule~\ref{ch03:rule:name-identity} in any of the three directions.
Similarly, the set-and-generator fixtures of
Checkclasses~\ref{ch03:check:objects} and~\ref{ch03:check:arrows} pin
Rule~\ref{ch03:rule:ordered-input} at every constructor that
Section~\ref{ch03:sec:ordered-rule} enumerated, and the boolean-index
fixture pins the subclass subtlety of Section~\ref{ch03:sec:arrows}.

\subsection{The validation surface in one table}

Table~\ref{ch03:tab:validation} summarizes the constructor validation
surface that the fixtures pin: each order-sensitive or type-sensitive
input position, what it rejects, and with which typed error. The table
is descriptive of the Gate~1 prototype sources; the authoritative
statement is the code bound by the manifest, and the fixtures are what
tie the two together.

\begin{table}[htbp]
\centering
\small
\begin{tabular}{@{}llll@{}}
\toprule
Constructor & Input & Rejected forms & Error \\
\midrule
\texttt{FiniteObject} & \texttt{name} & blank, non-string & \texttt{ValidationError} \\
\texttt{FiniteObject} & \texttt{elements} & set, mapping, iterator, string & \texttt{ValidationError} \\
\texttt{FiniteObject} & \texttt{elements} & duplicates, unhashables & \texttt{ValidationError} \\
\texttt{FiniteArrow} & \texttt{name} & blank, non-string & \texttt{ValidationError} \\
\texttt{FiniteArrow} & \texttt{source}, \texttt{target} & non-objects & \texttt{ValidationError} \\
\texttt{FiniteArrow} & \texttt{mapping} & unordered containers & \texttt{ValidationError} \\
\texttt{FiniteArrow} & \texttt{mapping} & wrong length (totality) & \texttt{ValidationError} \\
\texttt{FiniteArrow} & \texttt{mapping} & bool, non-int, out of range & \texttt{ValidationError} \\
\texttt{compose} & middle objects & any object inequality & \texttt{CompositionTypeError} \\
\texttt{Endomap} & \texttt{arrow} & source $\neq$ target & \texttt{CompositionTypeError} \\
\texttt{Endomap.power} & exponent & negative, bool, non-int & \texttt{ValidationError} \\
\texttt{SemigroupoidPath} & \texttt{arrows} & empty, non-arrows & \texttt{ValidationError} \\
\texttt{SemigroupoidPath} & \texttt{arrows} & adjacent mismatch & \texttt{CompositionTypeError} \\
\bottomrule
\end{tabular}
\caption{Constructor validation surface of \texttt{category.py}.
``Object inequality'' includes all three coercion variants: name-only,
carrier-only, and order-only differences.}
\label{ch03:tab:validation}
\end{table}

\subsection{What negative fixtures establish, and what they do not}

A negative fixture establishes that one representative of an invalid
class is rejected on the tested platform. It does not establish that
every member of the class is rejected --- that would require the
exhaustive or formal treatment that positive properties receive in
Chapters~10 and~11 --- and the suite's design acknowledges this by
choosing representatives at the \emph{boundary} of each class (length off
by exactly one; index exactly $|Y|$; the same-carrier coercion cases),
where an implementation bug is most likely to live. The honest summary,
maintained throughout the monograph's evidence chapters: positive
behavior on the bounded families is checked exhaustively; rejection
behavior is checked by boundary representatives; neither is
proof-assistant verification of the implementation, and the Lean seed's
verified scope (Chapter~11) covers the mathematical laws, not the Python
constructors. Within those limits, the fixtures make the chapter's three
rules --- name-inclusive identity, ordered input, names-out-of-equality
--- into properties that cannot regress silently: any future edit that
weakens a rule turns at least one fixture red.

\section{Summary}\label{ch03:sec:summary}

The chapter's content in one table of commitments:

\begin{itemize}
\item An object is a (name, ordered carrier) pair; equality includes both
(Definition~\ref{ch03:def:object}, Rule~\ref{ch03:rule:name-identity});
the argument is anti-coercion, with the register, refactor, and units
examples (Examples~\ref{ch03:ex:two-registers}--\ref{ch03:ex:units}) as
the canonical failure patterns.
\item Order-sensitive constructor inputs --- seven positions, enumerated
--- must be replayable ordered sequences; sets, mappings, one-shot
iterables, and strings are rejected at construction
(Rule~\ref{ch03:rule:ordered-input}); the constructed values are then
pure functions of their arguments (Invariant~\ref{ch03:inv:replay}).
\item Hash-order iteration is a general software-evidence threat
(Threat~\ref{ch03:threat:hash-order}): it converts digests from evidence
into noise, and the kernel's refusal-at-the-boundary is the only response
that keeps canonical forms independent of ambient interpreter state.
\item An arrow is its judgment --- declared source, declared target,
mapping tuple; equality is extensional over exactly this data and is a
congruence for composition and application
(Definitions~\ref{ch03:def:arrow}--\ref{ch03:def:ext-eq},
Proposition~\ref{ch03:prop:congruence}); presentation names are
diagnostics only and never enter equality
(Rule~\ref{ch03:rule:names-out}).
\item Every rule above is pinned by negative fixtures asserting exact
typed errors at exact boundaries
(Checkclasses~\ref{ch03:check:objects}--\ref{ch03:check:ops}), with the
evidential status \textsc{executable validation} and the limits of that
status stated.
\end{itemize}

With objects and arrows fixed and the determinism discipline in force,
Chapter~4 can state and prove the category laws for the kernel exactly as
Chapter~\ref{ch02:chap} proved them for bare finite functions --- the
typing adds hypotheses, never obstacles --- and Chapters~5 and~6 can
erect the endomap/path separation on ground that no longer shifts.
